\subsection{導入：ソフトウェア保守が重要な理由}
ソフトウェア開発が成功すると、利用者要求を満たすソフトウェア製品が提供される。しかし、運用が始まると、潜在していた欠陥が見つかり、運用環境が変化し、新しい利用者要求が現れる。したがって、ソフトウェア製品は進化しなければならない。\par
保守段階は、保証期間または導入後支援の終了後に始まると考えられがちである。しかし、実際の保守活動はもっと早く始まる。納品前から、運用後の支援、保守性、移行手順、支援体制を計画する必要がある。\par
近年は、開発、保守、運用の分断を減らし、継続的にソフトウェアを進化させる考え方が強くなっている。DevOps の実践やツールが普及したことで、開発者、運用担当者、保守担当者が協力し、素早く安全に変更を届けることが重視されている。\par
\begin{pointbox}{重要}保守は「古いソフトウェアを仕方なく直す作業」ではない。運用中のソフトウェアの価値を維持し、変化に対応し、利用者に継続的な便益を届けるための中核的なソフトウェア工学活動である。\end{pointbox}
